% ==================================================================
%  SECCIÓN 4.1 — Caracterización y preparación del corpus (EDA)
%  Fragmento listo para \input{} dentro del Capítulo 4 de la TESIS.
%
%  REQUISITOS DEL PREÁMBULO DE LA TESIS (ya presentes en el plan):
%    \usepackage{graphicx,booktabs,tabularx,xcolor,float,caption}
%    \definecolor{g0}{HTML}{17365D} \definecolor{g2}{HTML}{2AA8A8}
%    \definecolor{g3}{HTML}{38BF8F} \definecolor{g4}{HTML}{7FD16B}
%    \graphicspath{{figuras/}{../ENTREGA_CLASSROOM/Limpieza_Dataset/figuras/}}
%
%  Las figuras fig1..fig7 son las generadas por eda_limpieza.py.
%  Las cifras de normalización y partición provienen de particionar_dataset.py.
% ==================================================================

\section{Caracterización y preparación del corpus (EDA)}
\label{sec:eda}

La validez de cualquier métrica del sistema depende de la calidad y de la partición
del corpus con que se entrena y evalúa. Por ello, antes de la experimentación, se aplicó
un protocolo de análisis exploratorio de datos (EDA) de cinco pasos al subconjunto
\texttt{validated} de Common Voice 27.0, variedad quechua sureño (\texttt{qxp}), que
constituye el corpus base del módulo de reconocimiento de voz (ASR). El punto de partida
fueron \textbf{22\,727} clips validados; cada clip aporta su audio, su transcripción y
metadatos del hablante.

\subsection{Diagnóstico de calidad}
\label{sec:eda-calidad}

El examen de valores faltantes, duplicados y errores reveló cuatro hechos con
consecuencias directas para el entrenamiento:

\begin{itemize}\itemsep2pt
  \item \textbf{Columnas inservibles.} \texttt{variant} y \texttt{segment} están vacías
        al 100\,\%, y \texttt{accents} al 63\,\%; se descartan o se marcan como
        \emph{unknown} sin pérdida de clips útiles.
  \item \textbf{Desbalance severo de género.} 16\,079 clips \emph{female} frente a solo
        308 \emph{male} (Figura~\ref{fig:barras}). El modelo tenderá a reconocer mejor
        voces femeninas; es una amenaza de validez para un uso médico con pacientes de
        ambos sexos.
  \item \textbf{Concentración de hablantes.} Un solo hablante aporta el 8,9\,\% de los
        clips, lo que obliga a controlar la partición para que su voz no infle las
        métricas.
  \item \textbf{Reúso masivo de frases.} Las 22\,727 grabaciones corresponden a solo
        \textbf{2\,058} frases distintas (cada una grabada $\sim$11 veces). Esto no es
        un error, pero crea un riesgo de \emph{fuga de datos} que se trata en la
        Sección~\ref{sec:eda-particion}.
\end{itemize}

\begin{figure}[H]
  \centering
  \includegraphics[width=0.9\linewidth]{fig1_nulos.png}
  \caption{Mapa de valores faltantes por columna. El bloque de columnas vacías y el
  faltante demográfico orientan las decisiones de limpieza.}
  \label{fig:nulos}
\end{figure}

\begin{figure}[H]
  \centering
  \includegraphics[width=0.92\linewidth]{fig4_barras_categoricas.png}
  \caption{Distribución de variables categóricas. El desbalance de género y la
  concentración de hablantes son estructurales del corpus.}
  \label{fig:barras}
\end{figure}

\subsection{Anatomía numérica y correlaciones}
\label{sec:eda-correlaciones}

El análisis univariado (histogramas y diagramas de caja) sitúa la duración mediana en
$\sim$4,6\,s y la velocidad de habla en $\sim$7,2\,car/s, con colas que la regla
intercuartílica (IQR, criterio de Tukey) marca como atípicas. El análisis bivariado
contrasta los coeficientes de Pearson (lineal) y Spearman (monótona, robusta): mientras
\texttt{n\_chars}\,$\sim$\,\texttt{dur\_s} coinciden en 0,71, el par
\texttt{down\_votes}\,$\sim$\,\texttt{net\_votes} diverge ($-0{,}37$ frente a $-0{,}50$),
lo que justifica reportar ambas medidas y no fiarse de una sola
(Figura~\ref{fig:heatmap}).

\begin{figure}[H]
  \centering
  \includegraphics[width=0.96\linewidth]{fig6_heatmaps.png}
  \caption{Mapas de calor de correlación de Pearson y de Spearman. La discrepancia entre
  ambos advierte contra interpretar relaciones no lineales con un único coeficiente.}
  \label{fig:heatmap}
\end{figure}

\subsection{Variable objetivo: dominio de salud}
\label{sec:eda-salud}

Como la tesis se orienta a la consulta médica, se analizó el subconjunto de dominio
\emph{healthcare}, que cubre el \textbf{77,8\,\%} del corpus limpio: una cobertura
favorable. Las pruebas de Mann--Whitney detectan diferencias significativas, aunque
pequeñas, entre los clips de salud y el resto en duración y longitud
($p<10^{-30}$), coherentes con un vocabulario clínico algo más breve.

\subsection{Limpieza aplicada}
\label{sec:eda-limpieza}

Sobre la evidencia anterior se aplicó una limpieza conservadora: se eliminaron columnas
vacías, constantes y redundantes; se normalizó la etiqueta multi-valor
\texttt{sentence\_domain} (corrigiendo 76 filas con tokens duplicados) hacia los
indicadores \texttt{domains}/\texttt{is\_healthcare}; los faltantes demográficos se
marcaron como \emph{unknown} en lugar de borrar la fila (hacerlo habría costado hasta el
63\,\% del corpus); y se excluyeron los clips atípicos por IQR en \texttt{dur\_s} y
\texttt{cps} junto con los de \texttt{net\_votes}\,$\le$\,0. El corpus pasó de
\textbf{22\,727} a \textbf{21\,653} clips (Tabla~\ref{tab:limpieza}).

\begin{table}[H]
  \centering
  \small
  \caption{Efecto de la limpieza sobre el tamaño del corpus.}
  \label{tab:limpieza}
  \begin{tabular}{l r r}
    \toprule
    Etapa & Clips & Variación \\
    \midrule
    Validado (crudo)            & 22\,727 & --- \\
    Tras excluir atípicos (IQR) & 21\,653 & $-1\,074$ ($-4{,}7$\,\%) \\
    \bottomrule
  \end{tabular}
\end{table}

\subsection{Normalización ortográfica del quechua}
\label{sec:eda-normalizacion}

El WER/CER de un ASR es sensible a diferencias ortográficas invisibles en la referencia.
Se fijó una convención única y documentada: forma Unicode NFC; unificación de las
variantes tipográficas del apóstrofo al apóstrofo recto \texttt{'} —preservando su valor
\emph{fonémico} como marca de glotalización en las oclusivas quechuas (\emph{p'}, \emph{t'},
\emph{ch'}, \emph{k'}, \emph{q'})—; conversión a minúsculas; retiro de la puntuación no
pronunciada conservando apóstrofo y guion internos; y colapso de espacios. La
normalización afectó al \textbf{99,9\,\%} de las transcripciones (casi todas por el paso a
minúsculas o la puntuación) y redujo las frases únicas de 2\,062 a \textbf{2\,058} al
fundir variantes equivalentes. No se hallaron dígitos en el texto, por lo que no fue
necesario desambiguar números hablados.

\subsection{Partición anti-fuga: el \emph{trade-off} hablante\,/\,frase}
\label{sec:eda-particion}

El reúso masivo de frases obliga a controlar la \emph{fuga de datos}. Se evaluaron los dos
ejes de partición posibles, cada uno disjunto en un eje y midiendo la fuga residual en el
otro (Tabla~\ref{tab:particion}). El hallazgo es que ambos ejes son
\textbf{incompatibles} en este corpus: como casi todos los hablantes graban las mismas
2\,058 frases, una partición disjunta por hablante deja el $\sim$100\,\% de las frases de
prueba también en entrenamiento, y una disjunta por frase deja el $\sim$100\,\% de las
voces compartidas.

\begin{table}[H]
  \centering
  \small
  \caption{Las dos particiones (80/10/10 por número de clips) y su fuga residual.}
  \label{tab:particion}
  \begin{tabularx}{\linewidth}{@{}l r r r X@{}}
    \toprule
    Partición & Train & Val & Test & Fuga residual (test) \\
    \midrule
    Disjunta por \textbf{hablante} & 17\,139 & 2\,260 & 2\,254 & 100\,\% de las \emph{frases} ya vistas \\
    Disjunta por \textbf{frase}    & 17\,309 & 2\,168 & 2\,176 & 100\,\% de las \emph{voces} ya vistas \\
    \bottomrule
  \end{tabularx}
\end{table}

Para un ASR médico, cuyo objetivo es transcribir \emph{frases nuevas} dichas en consulta,
la amenaza más grave es que el modelo memorice el repertorio fijo de frases. Por ello se
adopta la \textbf{partición disjunta por frase} como principal, y se conserva la disjunta
por hablante como evaluación complementaria de robustez ante voces no vistas; el
solapamiento de hablantes se declara como limitación. Ambas particiones quedan fijadas en
columnas reproducibles del dataset (\texttt{split\_sentence}, \texttt{split\_speaker}),
con semilla \texttt{SEED=42}.

\subsection{Síntesis e implicaciones para la experimentación}
\label{sec:eda-sintesis}

El EDA entrega un corpus de 21\,653 clips depurados, con texto normalizado y dos
particiones verificadas sin fuga en su eje respectivo. De él se derivan tres decisiones
que gobiernan el resto del Capítulo~4: (i) evaluar con la partición disjunta por frase y
reportar también la disjunta por hablante; (ii) informar el WER \emph{estratificado por
género} para exponer el sesgo derivado del desbalance 16\,079\,/\,308; y (iii) usar este
mismo corpus depurado —no el crudo— como base del factor «calidad del corpus» en el
experimento de ablación (Sección~\ref{sec:ablacion} de la metodología). Todo el proceso es
reproducible con \texttt{eda\_limpieza.py} y \texttt{particionar\_dataset.py}.
